\chapter{Fabrication Spaces, Resources, and Manufacturing}

\section{Labriola Innovation Hub Overview}
The Labriola Innovation Hub provides dedicated spaces for capstone teams, both for project work during class time and for extended hours outside of class. You must complete required training before accessing any shop or workspace.

\textbf{Thorson Capstone Design Center}: primary workspace for all capstone teams, reserved for capstone students only. This open design space doubles as your build space and after-hours workspace. Blastercard access required. See the Labriola Canvas page for hours, Grad TA contacts, locker assignments, and storage policies. The Thorson Center is available during posted hours with TA supervision.

\textbf{xWorks Project Bays}: dedicated bays for teams working on vehicles or large-footprint projects (e.g., concrete canoe, Formula SAE car). Each qualifying team is assigned one bay for the semester. Review the expectations page for Blastercard enrollment.

\textbf{Labriola Conference Rooms}: second-floor conference rooms reserved for capstone students TTH 8:00--11:00 AM for Sprint Reviews and Design Reviews only. See the Labriola page for reservation procedures and limitations.

\section{InnoHub Shops and Makerspace}
The InnoHub includes specialized shops, each staffed and equipped for specific fabrication processes:

\textbf{Metal Shop}: manual mills, lathes, drill presses, band saws, grinders, MIG and TIG welders, sheet metal equipment. CNC mills and lathes available after completing manual machine training first.

\textbf{Wood Shop}: table saw, band saw, wood lathe, drill press, sanders, routers. CNC router available with training.

\textbf{Electronics Shop}: soldering stations, hot air rework station, oscilloscopes, bench power supplies, component storage, PCB assembly workspace.

\textbf{Composites Shop}: vacuum bagging, wet layup, curing oven. Specialized safety training required due to resin sensitization risks.

\textbf{Makerspace}: FDM and resin 3D printers, laser cutters, vinyl cutters, general hand tools, and fabrication workspace.

\begin{warningbox}[Training Is Mandatory; No Exceptions]
Complete machine-specific safety training on the \textbf{Labriola InnoHub Canvas site} for every piece of equipment you plan to use. No drop-ins are accepted; you will be turned away. Trainings from other facilities (ME Machine Shop, physics lab, personal experience) do \textbf{not} substitute for InnoHub training. CNC training requires completion of the corresponding manual machine training first. Training prioritization: if a TA needs a machine or work area for safety training, you will be asked to return at a different time.
\end{warningbox}

\textbf{Extended machine use}: if you plan to use manual mills/lathes, welders, wood lathe, or table saw for more than a \textbf{4-hour period}, get manager approval in advance.

\section{Manufacturing Review Meeting}
If your team plans to build your final project in \emph{any} Labriola facility (Metal Shop, Wood Shop, Electronics Shop, Composites Shop, or Makerspace) you \textbf{must schedule a 30-minute Manufacturing Review Meeting} with one of the shop managers or the operations manager before ordering materials.

This meeting is a \textbf{required checkoff item on your purchasing form}. It must be completed before CP will process material orders for your build. Bring your design materials (CAD, drawings, BOM) to the meeting. The InnoHub team will:
\begin{itemize}[nosep,leftmargin=1.5em]
\item Provide feedback and guidance on your manufacturing procedures
\item Suggest additional tools you may need
\item Recommend alternative materials or approaches that may save time and money
\item Identify which InnoHub machines and resources can improve your process
\end{itemize}

Schedule through the booking link on Canvas. The InnoHub team has blocked time specifically for these meetings. Do not wait until the last minute; scheduling fills up quickly around CDR.

\section{3D Printing Policies}
3D printing is the most commonly used fabrication method in capstone. InnoHub policies balance access across all teams:
\begin{itemize}[nosep,leftmargin=1.5em]
\item Each team may use \textbf{one printer at a time}.
\item Teams receive \textbf{1,000\,g of PLA Basic filament free}. Usage subject to standard time and gram limits.
\item Individual prints must be $\leq$\textbf{250\,g} without TA or manager approval.
\item \textbf{Daytime prints}: up to 5 hours unless approved for exception.
\item \textbf{Overnight prints}: up to 12 hours, prior TA approval required.
\item After exhausting the free allotment: purchase filament rolls or pay per print. Purchased filament still subject to time limits.
\item \textbf{Advanced materials}: TPU (flexible) and PETG (chemical resistant) available at cost, as are resin printers (SLA/DLP) for high-detail parts.
\item \textbf{Payment}: all payments processed \emph{before printing} via worktag transfer or credit card through the InnoHub payment portal.
\end{itemize}

\section{Paint Booth and Surface Finishing}
Spray painting must use the \textbf{designated paint booth} only. Complete the paint booth pre-quiz on Canvas before your first use. Sign up for work time through the booking system.

\section{ME Machine Shop}
The ME Machine Shop and Plasma Cutter in Brown Hall have \textbf{limited access} due to ME course usage. Use InnoHub as your primary fabrication resource.

\section{General InnoHub Policies}
\begin{itemize}[nosep,leftmargin=1.5em]
\item Shops available during posted hours with TA supervision only.
\item \textbf{TA authority}: Labriola TAs have authority in all shops. If they provide guidance or request something, comply promptly. Contact Julia Roos (jroos@mines.edu), InnoHub Operations Manager, for policy concerns.
\item \textbf{Tool Crib}: return all tools by closing time daily. Purchase fasteners from project budget after prototyping.
\item \textbf{Cleanliness}: clean workspace before leaving. No overnight storage without approval.
\item Some materials and services available for purchase from InnoHub directly.
\end{itemize}


\section{Manufacturing Process Selection}
Choosing the right manufacturing process for each component of your design is a key engineering decision. Consider:

\textbf{3D Printing (FDM)}: best for complex geometries, rapid iteration, low-volume production. Lead time: hours to days. Limitations: anisotropic strength (weak between layers), limited material properties, surface finish requires post-processing. Use for: enclosures, brackets, jigs, fixtures, form-fit prototypes.

\textbf{Laser cutting}: best for flat parts from sheet material (acrylic, plywood, MDF, thin metal). Lead time: minutes to hours. Excellent dimensional accuracy in-plane ($\pm$0.1\,mm). Use for: panels, structural frames, gaskets, decorative elements.

\textbf{CNC machining}: best for metal parts requiring high strength, tight tolerances, and good surface finish. Lead time: hours to days (InnoHub) or weeks (outsourced). Use for: structural brackets, motor mounts, shaft couplings, bearing housings, heat sinks.

\textbf{Hand fabrication}: sheet metal bending, drilling, tapping, filing, soldering, adhesive bonding. Often the fastest approach for one-off parts. Lead time: minutes to hours. Use for: simple brackets, wire harnesses, electrical enclosures, quick modifications.

\textbf{Outsourced fabrication}: PCB manufacturing (JLCPCB, PCBWay), sheet metal fabrication (SendCutSend, OSH Cut), CNC machining (Xometry, Protolabs). Lead time: 1--4 weeks including shipping. Use when the required process or material is not available in the InnoHub or when volume justifies the cost.

\begin{tipbox}
When choosing a manufacturing process, consider the \textbf{whole lifecycle}: not just the first prototype, but also replacement parts, maintenance, and potential future production. If the client may eventually manufacture your design at scale, design for the production process (injection molding, die casting) even if the prototype uses a different method (3D printing, machining).
\end{tipbox}

\section{Quality Control During Fabrication}
Professional fabrication includes quality checkpoints:

\textbf{Incoming inspection}: when components arrive, verify: correct part number, correct quantity, no shipping damage, and visual quality. For precision parts, measure critical dimensions with calipers before assembly. A bracket that is 0.5\,mm too thick will not fit the enclosure.

\textbf{In-process checks}: verify dimensions after each machining operation. Check 3D prints for warping, layer adhesion, and dimensional accuracy before assembling. Test electrical sub-assemblies before integration.

\textbf{Final inspection}: before the system is considered ``built,'' inspect all fasteners for proper torque, all electrical connections for continuity, all fluid connections for leaks, and all moving parts for proper clearance and alignment.

Document quality issues and their resolution. At FDR, the as-designed vs. as-built comparison (Chapter~10) should include any fabrication deviations and their impact on performance.



\begin{tipbox}[Student Guide Cross-Reference]
For the InnoHub training sequence and Manufacturing Review Meeting timeline, see the \textbf{Student Guide, SDI Sprint~4} and \textbf{SDII Sprint~8} sections.
\end{tipbox}

\section{Your First Week in the InnoHub}
If this is your first time using the InnoHub shops, follow this sequence:

\begin{enumerate}[nosep,leftmargin=1.5em]
\item \textbf{Enroll} in the Labriola InnoHub Canvas site (link on the main Capstone Canvas page).
\item \textbf{Complete general safety training} on the Canvas site.
\item \textbf{Complete machine-specific training} for every piece of equipment you plan to use. Each machine has its own training module. Allow 1--2 weeks for training availability.
\item \textbf{Schedule your Manufacturing Review Meeting} with a shop manager (required before ordering build materials).
\item \textbf{Get your Blastercard activated} for Thorson Design Center access.
\item \textbf{Locate your team's locker/storage} in the Thorson Center.
\item \textbf{Introduce yourself to the Grad TAs}; they are your primary support for tools, safety questions, and fabrication guidance.
\end{enumerate}

\textbf{TA contact information}: the current semester's Capstone Design Grad TAs are listed on the Labriola Spaces page on Canvas. These TAs support teams with hand and power tool needs, safety questions, and storage. If TAs cannot assist you, contact Prof.\ Kristy Csavina (Capstone Design Director, kcsavina@mines.edu). For general InnoHub questions, contact innohub@mines.edu.

\section{Fabrication Tips from Past Teams}
\begin{itemize}[nosep,leftmargin=1.5em]
\item \textbf{Prototype early, prototype often}: your first 3D print will reveal design problems that CAD did not. Print a rough version early and iterate.
\item \textbf{Test-fit before committing}: before machining the final aluminum bracket, make a cardboard or PLA test piece to verify dimensions, clearances, and mounting hole alignment.
\item \textbf{Label everything}: label every wire, every connector, every tube, every PCB. During debugging at 11\,PM, clear labels save hours.
\item \textbf{Take photos during assembly}: document the build process with photos at every stage. These are essential for the FDR as-built documentation and the O\&M manual.
\item \textbf{Clean as you go}: a messy workspace leads to lost parts, damaged components, and safety hazards. Spend the last 10 minutes of every build session organizing your workspace.
\end{itemize}



\section{Outsourced Fabrication Services}
When InnoHub cannot provide the required process, or when the quality requirements exceed student fabrication capabilities, outsourced fabrication is an option:

\textbf{PCB fabrication}: JLCPCB (cheapest for prototypes: 5 boards for \$2--5 + shipping), PCBWay (more options, good for assembly), OSH Park (US-based, higher quality, \$5/in$^2$ for 3 boards). All accept Gerber files generated from KiCad, Altium, or EasyEDA.

\textbf{Sheet metal}: SendCutSend (laser cutting and bending from uploaded DXF files, 1--2 week turnaround), OSH Cut (similar), and local sheet metal shops for complex assemblies.

\textbf{CNC machining}: Xometry and Protolabs offer instant quoting from CAD files. More expensive than InnoHub but higher precision and faster turnaround. Use for parts that require tolerances tighter than $\pm$0.05\,mm or materials not available in the InnoHub.

\textbf{Waterjet and laser cutting}: available through SendCutSend and local vendors. Excellent for thick metal plates, complex profiles, and materials that cannot be machined conventionally (composites, ceramics).

When outsourcing, add the vendor lead time to your procurement timeline and budget the fabrication cost explicitly. Order from outsourced fabrication vendors at CDR at the latest to have parts in hand for Sprint~10 assembly.

\section{Working Safely with Power Tools}
Even with completed InnoHub training, maintain constant vigilance with power tools:

\textbf{The 5-second rule}: before starting any power tool, pause for 5 seconds and mentally confirm: (1)~workpiece is secure, (2)~tool guards are in place, (3)~PPE is on, (4)~path of tool/workpiece is clear of hands and obstructions, (5)~emergency stop is accessible. This mental checklist prevents the most common shop injuries.

\textbf{Never reach across a running tool}: turn off the tool and wait for it to stop completely before adjusting the workpiece, clearing chips, or measuring.

\textbf{Fatigue awareness}: tired people make mistakes. If you have been in the shop for more than 4 hours, take a break. If you are working late at night, consider whether the task can wait until morning when you are alert. The most serious shop injuries occur during the last hour of long sessions.



\section{CNC Machining Tips for Capstone}
If your project involves CNC-machined parts (aluminum brackets, motor mounts, bearing housings), these practices save time and money:

\textbf{Design for 3-axis}: unless you have specific 4- or 5-axis requirements, design parts that can be machined in a single setup on a 3-axis mill. This means: features accessible from one direction, no undercuts requiring repositioning, and flat reference surfaces for vise clamping.

\textbf{Stock material selection}: start with near-net-shape stock to minimize material removal. A bracket cut from 0.5'' plate requires much less machining than one cut from a 2'' block. Standard stock sizes: plate (various thicknesses), bar (round, square, hex), and angle/channel (structural shapes).

\textbf{Coordinate with InnoHub staff}: discuss your part with a shop manager before programming the CNC. They can suggest: optimal tool paths, appropriate feeds and speeds for the material, fixturing strategies, and whether the part is better made on the CNC or by manual machining (sometimes a simple bracket is faster on the manual mill).

\textbf{Bring drawings}: even though the CNC runs from G-code, bring a dimensioned drawing to the shop. The machinist (you or the TA) needs to verify the setup, check the first article, and inspect the finished part against the drawing.

\section{Welding Considerations}
Welding joins metal permanently with high strength. InnoHub offers MIG and TIG welding with training:

\textbf{MIG welding}: easiest to learn, fastest process. Good for steel and aluminum $\geq$1/8'' thick. Produces adequate structural welds for most capstone applications. Not suitable for thin sheet metal ($<$1/16'') or high-precision cosmetic applications.

\textbf{TIG welding}: slower, requires more skill, produces cleaner and more controlled welds. Necessary for thin material, stainless steel, and applications where weld appearance matters. Requires more training time than MIG.

\textbf{Design for welding}: provide adequate access for the welding torch (the nozzle is $\sim$1'' diameter). Specify weld type (fillet, butt, corner) and size on your drawings. Design joints with adequate material thickness (minimum 1/8'' for structural MIG welds). Account for heat distortion: weld assemblies may warp; design fixtures to hold parts in alignment during cooling.

\textbf{Safety}: welding produces intense UV radiation (arc eye), hot spatter (burns), and toxic fumes (metal oxide particles). Full PPE required: welding helmet with auto-darkening lens ($\geq$shade 10), welding gloves, long sleeves (leather jacket recommended), closed-toe boots, and adequate ventilation. Never weld near flammable materials, compressed gas cylinders, or battery packs.



\section{3D Print Design Optimization for Capstone}
Beyond the basic guidelines in Chapter~4, optimize your 3D prints for capstone-specific constraints:

\textbf{Print orientation for function}: orient parts so that: (1)~the primary load is perpendicular to layer lines (strongest direction), (2)~the most important surface faces up or away from the build plate (best surface finish), (3)~support material is minimized (less post-processing, less waste of your 1,000\,g free allotment).

\textbf{Infill strategy}: 20\% infill is the default; adequate for most non-structural parts. For structural parts (brackets, mounts, enclosures under stress): use 50--100\% infill with 4+ perimeter walls. For weight-sensitive parts: use 10\% infill with gyroid pattern. For water-tight parts: use 100\% infill or vapor-smooth the exterior.

\textbf{Iteration strategy}: print version 1 at 0.3\,mm layer height (fast, rough) to check fit and form. Print version 2 at 0.2\,mm (standard) after confirming dimensions. Print the final version at 0.15\,mm (fine) only if surface quality matters for the client or showcase. Each quality step doubles print time; use the lowest quality that meets your need.

\textbf{Heat-set inserts}: for any part that will be assembled/disassembled more than once, use heat-set brass inserts (M3 are the most common in capstone). Press them into undersized holes using a soldering iron at 220\,\degC{} for PLA, 250\,\degC{} for PETG. The result is a metal thread in a plastic part that survives dozens of assembly cycles.

\textbf{Combining 3D printing with other processes}: print the complex geometry (organic shapes, internal channels, mounting features) and machine the precision interfaces (bearing seats, shaft bores, sealing surfaces). This hybrid approach leverages the strengths of each process.

\section{Laser Cutting for Capstone}
Laser cutters in the InnoHub Makerspace cut flat stock quickly and accurately:

\textbf{Materials}: acrylic (3--6\,mm), plywood (3--6\,mm), MDF, cardboard, felt, fabric. The laser cutter cannot cut metal (except thin aluminum foil on some machines) or polycarbonate (produces toxic fumes).

\textbf{Design for laser}: use 2D CAD (DXF or SVG format). All geometry must be in a single plane. Design with kerf compensation: the laser removes $\sim$0.1--0.2\,mm of material width. For press-fit joints, offset the slot by half the kerf width.

\textbf{Tab-and-slot construction}: laser-cut panels can be assembled into 3D structures using tab-and-slot joints. Design tabs with a slight interference fit (0.05--0.1\,mm). Glue with cyanoacrylate (instant bond) or acrylic cement (chemical weld for acrylic-to-acrylic joints). This technique produces rigid enclosures and frames in a fraction of the time of 3D printing.

\textbf{Living hinges}: thin cuts in plywood or acrylic create flexible hinges. Useful for enclosure lids and folding structures. Design hinge width based on material: 0.5\,mm for 3\,mm acrylic (fragile), 1\,mm for 3\,mm plywood (more durable).



\section{Electronics Assembly in the InnoHub}
The Electronics Shop in the InnoHub provides soldering stations, a hot air rework station, and component storage. Best practices for electronics assembly:

\textbf{Soldering technique}: use a temperature-controlled iron at 300--350\,\degC{} for leaded solder or 350--380\,\degC{} for lead-free. Apply flux to the joint, heat the pad and component lead simultaneously (not just one), then touch solder to the heated joint (not to the iron). A good joint is shiny (leaded) or slightly matte (lead-free), concave, and wets both the pad and the lead.

\textbf{Common soldering defects}: cold joint (dull, grainy (insufficient heat), solder bridge (excess solder connecting adjacent pads) use solder wick to remove), tombstoned component (one end lifted (reheat both pads), insufficient solder (pad visible through the solder) add more). Inspect every joint under 10$\times$ magnification before powering the board.

\textbf{SMD soldering by hand}: for SOT-23, SOIC, and 0805/0603 passives, use fine-tip solder and a magnifier. Tack one pad first, then solder the remaining pads. For QFP and TQFP packages ($\leq$0.5\,mm pitch), use drag soldering with flux: apply generous flux, tack one corner pin, then drag the solder-loaded iron tip across the row of pins. The flux prevents bridges; use solder wick to clean any that form.

\textbf{Through-hole assembly}: solder from the bottom of the board. Insert the component from the top, flip the board, solder, and trim the leads. For multi-pin connectors, tack one pin first to verify alignment before soldering the rest.

\textbf{ESD protection}: electrostatic discharge can damage CMOS ICs invisibly; the component appears functional but has reduced reliability. Use an ESD wrist strap connected to a grounded surface when handling bare ICs, PCBs, and sensitive modules. The InnoHub Electronics Shop provides ESD mats and wrist straps.



\section{Material Sourcing and Stock Preparation}
Before starting fabrication, ensure all materials are on hand and prepared:

\textbf{Metal stock}: order from McMaster-Carr or Online Metals. Allow 3--5 days for delivery. Verify dimensions with calipers upon receipt; raw stock can be slightly over or under nominal dimensions. Cut stock to rough size (leaving 3--5\,mm excess per side) before machining to final dimensions.

\textbf{3D printing filament}: verify you have enough filament remaining in your free allotment (check with InnoHub TAs). If additional filament is needed, purchase through the InnoHub payment portal before printing. Different colors of the same material have identical mechanical properties; choose color for aesthetics, not performance.

\textbf{Fastener kits}: order assortment kits (M3, M4, M5 metric, or \#4-40, \#6-32, \#8-32 imperial) rather than individual sizes. The per-unit cost is lower and you will inevitably need a size you did not anticipate. McMaster-Carr assortment kits are organized by size and stored in labeled compartments.

\textbf{Wire and cable}: for electronics projects, order hookup wire in at least 4 colors (red=VCC, black=GND, plus 2 signal colors), in both solid core (for breadboards) and stranded (for connections that flex). 22~AWG is the standard for signal wiring; use heavier gauge (18--14~AWG) for power distribution.

\section{Workspace Organization and Waste Management}
\textbf{The clean workspace principle}: a messy workspace is a slow workspace. You waste time searching for tools, risk losing small components, and increase the chance of assembly errors (wrong screw in the wrong hole because everything is jumbled together). Spend 10 minutes at the end of every build session organizing your workspace.

\textbf{Component organization}: use labeled bins, bags, or compartmentalized boxes for: received components (organized by vendor/order), installed components (keep a running installation checklist), spare components (clearly marked as spares), and rejected/damaged components (do not discard (they may be needed for failure analysis at FDR).

\textbf{Waste disposal}: dispose of materials properly. Metal chips and shavings: sweep into the designated metal waste bin (never vacuum) metal chips damage vacuum motors). 3D printing supports and failed prints: PLA is not recyclable through standard recycling; dispose in general waste unless your facility has a filament recycling program. Chemical waste (solvents, adhesives, resin): follow the SDS disposal instructions and use the designated chemical waste containers in the InnoHub.

\textbf{Tool return}: return all InnoHub tools to the Tool Crib before leaving each session. Check out tools through the designated system. If a tool is damaged during use, report it to the TA immediately; do not return it silently. The next team to use a damaged tool may be injured.



\section{The Build Phase Daily Workflow}
During active fabrication (Sprints 9--11), establish a daily workflow:

\textbf{Start of session} (10 min): review the task board. Identify today's build tasks and their dependencies. Gather required tools, materials, and drawings. Verify that all components for today's assembly are available.

\textbf{Build activities} (2--4 hours): execute the build tasks. Follow the assembly sequence documented in the manufacturing plan. Take photographs at each major assembly step. Record any deviations from the plan (``drilled hole 2\,mm left of drawing dimension to avoid wire interference'').

\textbf{Test during build}: after each subsystem is assembled, perform a basic functionality check before proceeding to the next subsystem. Catching a wiring error after assembling one subsystem is a 30-minute fix; catching it after the enclosure is sealed is a 3-hour disassembly.

\textbf{End of session} (15 min): update the task board. Log what was completed, what remains, and any blockers. Clean the workspace. Return tools. Store the project securely.

\textbf{Documentation checkpoint} (weekly): at the end of each week, update: the as-built documentation (photos, deviation notes), the BOM (mark installed components), and the test log (subsystem test results). This weekly checkpoint prevents documentation debt from accumulating to unmanageable levels before FDR.


\section{Practice Questions}
\begin{enumerate}[leftmargin=1.5em]
\item List the training sequence required before you can use the CNC mill in the InnoHub Metal Shop.
\item You have used 800\,g of your free PLA allotment and need to print a 400\,g enclosure. What are your options and required steps?
\item Your Manufacturing Review revealed that your planned sheet metal bending is not feasible on InnoHub equipment. Propose three alternative fabrication approaches.
\end{enumerate}


